\lhead{Cap\'itulo 1}
\rhead{\newtitle}
\cfoot{\thepage}
\renewcommand{\headrulewidth}{1pt}
\renewcommand{\footrulewidth}{1pt}
\chapter{Investigación IA}

\section{Introducción}
% La introducción es la presentación clara, breve y precisa del contenido general del proyecto del Electivo de Profundización I, que invita al lector a leer el documento de forma motivante, clara y precisa. Los aspectos que debe considerar son:
%    Las motivaciones para la elección del tema.
%    Antecedentes de la problemática o necesidad abordada con el proyecto.
%    El objetivo del proyecto mencionado al final de la introducción.
%    Todas las referencias necesarias para fundamentar y validar el tema expuesto.
\subsection{Contexto del problema}

% Migración en Chile
% Mercado laboral

\subsection{Problema de investigación}

% Discrepancia ideal vs real

\subsection{Justificación}

% Uso de ENE + CASEN
% Relevancia social y analítica



% \begin{center}
	%     \begin{minipage}{400px}
		%         Para cita en bloque
	%     \end{minipage}
% \end{center}

% \citep{CASEN2022, Castro-Romero2024}. 

\section{Objetivos}

% En esta parte se debe escribir el objetivo general y los objetivos específicos del tema a desarrollar.

\subsection{Objetivo General}

% Se debe declarar el objetivo general de todo el proyecto de manera que contenga de manera global a todos los objetivos específicos.



\subsection{Objetivos Específicos}

% Los objetivos específicos se deben desprender lógicamente desde el objetivo general, y corresponde a tareas específicas que conllevan a resolver el objetivo del proyecto.

% (Importante: que cada uno tenga vínculo directo con análisis/modelos)

\section{Hipótesis de Investigación}



\subsection{Hipótesis principal}



\subsection{Hipótesis secundarias}



\section{Metodología} 

% En la parte de metodología se deben explicar las herramientas técnicas, modelos matemáticos, análisis estadísticos que se utilizarán para resolver el problema, detallando:

%    Breve introducción a los modelos matemáticos, estadísticos o algorítmicos que se utilizarán.
%    Definición de variables, datos o modelos a utilizar.
%    Todas las referencias necesarias para fundamentar y sustentar teóricamente.


\subsection{Diseño de investigación}

% Tipo: cuantitativo, observacional, transversal

\subsection{Fuentes de datos}



\subsubsection{Encuesta Nacional de Empleo (ENE)}


\subsubsection{Encuesta de Caracterización Socioeconómica Nacional(CASEN)}

% Aquí justificas por qué usas ambas

\subsection{Definición de variables}



\subsubsection{Variable dependiente (y)}



\subsubsection{Variables independientes (X)}



\subsubsection{Variables de control}



\subsection{Preprocesamiento de datos}

% Incluye:

% limpieza de datos
% tratamiento de valores nulos
% transformación de variables
% generación de nuevas variables (feature engineering)

\subsection{Pipeline de análisis}



\subsection{Estrategia analítica}



\subsubsection{Análisis descriptivo}



\subsubsection{Modelos supervisados}



\subsubsection{Modelos no supervisados}


\section{Modelo computacional y algoritmo}

\subsection{Selección de modelos}

% \begin{lstlisting}[style=mystyle, caption={Modelo RandomForestClassifier en CASEN 2022}, label={01py}]
% 	# Mejor modelo encontrado para CASEN 2022 (Accuracy 0.68)
		
% 	rf = RandomForestClassifier(n_estimators=32, # numero de arboles
% 								criterion="gini", # medida para evaluar
% 								max_depth=8, # profundidad
% 								n_jobs=-1,
% 								verbose=1,
% 								random_state=20230728)
% \end{lstlisting}

\subsubsection{Regresión lineal}



\subsubsection{Regresión logística}



\subsubsection{K-Means}


\subsection{Justificación de modelos}

% por qué ese modelo
% qué problema resuelve

\subsection{Comparación de modelos}

% modelo considerado vs descartado
% justificación técnica

\subsection{Métricas de evaluación}

% MSE
% Accuracy
% Matriz de confusión
% Método del codo

\section{Resultados}



\subsection{Análisis descriptivo}

% tasas de ocupación
% ingresos
% gráficos

\subsection{Resultados de modelos supervisados}

% coeficientes
% probabilidades
% interpretación

\subsection{Resultados de clustering}

% perfiles laborales
% análisis de clusters

\subsection{Síntesis de resultados}

% conectar resultados entre sí
% no dejarlo fragmentado

\section{Conclusiones}



\subsection{Validación de hipótesis}

% confirmación o rechazo

\subsection{Implicancias}

% sociales
% económicas



\subsection{Limitaciones del estudio}

% datos observacionales
% sesgos
% falta de causalidad

\subsection{Líneas futuras}

% qué se podría mejorar